% ID: SEL-D0044-OB-E04
% ESTADO: cotejada. Segunda pasada visual y borrador compilado revisados.
% Fuente: Selectividad/sel_2014_matematicas_aplicadas/EXAMEN -ANDAL-3.pdf
% SHA-256: feb2acf9edc82ad3efbac4f8abf58b3e863451adddf04a82870066e873ab63fd
% Página: 2; opción B; ejercicio original 4.

\begin{enumerate}[label=\alph*)]
\item \textbf{(1.5 puntos)} Determine todas las muestras de tamaño 2 que, mediante muestreo aleatorio simple, se pueden extraer del conjunto $\{6,9,12\}$ y calcule la varianza de las medias de estas muestras.
\item \textbf{(1 punto)} Una empresa fabrica cuatro productos A, B, C y D, de los que elabora diariamente 40, 15, 25 y 120 unidades respectivamente.

Si un día se quiere elaborar una muestra de 40 unidades con los productos fabricados, por muestreo aleatorio estratificado con afijación proporcional, \textquestiondown{}qué número de unidades de cada producto se debe elegir?
\end{enumerate}
